\section{The finite integration technique}
\label{sec:theory}

\noindent\emph{Status: draft -- everything stated here is implemented and unit-tested; the derivation, the
motivating prose and the discussion of the non-orthogonal case are still being written.}

This section fixes the objects the implementation manipulates. It defines the two cell complexes of the grid, the
degrees of freedom as integrals over their elements, the three topological operators and the two identities they
satisfy exactly, the Maxwell grid equations, the diagonal Hodge stars that carry the material, and the way
boundary conditions and sources enter. The presentation follows Weiland \cite{weiland1977,weiland1996} and
\cite{clemens2001}; the energy and stability analysis that makes use of these objects is deferred to
the section on energy conservation and stability.

\subsection{The two cell complexes}

A structured hexahedral grid is the tensor product of three strictly increasing sequences of node coordinates,
$x_0 < \dots < x_{n_x}$, $y_0 < \dots < y_{n_y}$, $z_0 < \dots < z_{n_z}$. Its \emph{primal} objects are the
familiar ones: nodes at the corners, edges the segments joining adjacent nodes, faces the rectangles spanned by
two edges, cells the boxes spanned by three. Writing $n_k$ for the number of \emph{cells} along axis $k$, so that
axis $k$ carries $n_k+1$ nodes, the counts are those in Table~\ref{tab:counts}.

\begin{table}[H]
  \centering
  \caption{Entities of the primal and dual complexes of an $n_x\times n_y\times n_z$ grid, and their values for
  the $10\times7\times5$ cavity used in the validation runs.}
  \label{tab:counts}
  \begin{tabular}{@{}lrr@{}}
    \toprule
    Entity & Count & $10\times7\times5$ \\
    \midrule
    primal nodes   & $(n_x+1)(n_y+1)(n_z+1)$ & 528  \\
    primal edges   & $n_x(n_y+1)(n_z+1) + (n_x+1)n_y(n_z+1) + (n_x+1)(n_y+1)n_z$ & 1382 \\
    \quad of which $x$ & $n_x(n_y+1)(n_z+1)$ & 480 \\
    \quad of which $y$ & $(n_x+1)n_y(n_z+1)$ & 462 \\
    \quad of which $z$ & $(n_x+1)(n_y+1)n_z$ & 440 \\
    primal faces   & $(n_x+1)n_yn_z + n_x(n_y+1)n_z + n_xn_y(n_z+1)$ & 1205 \\
    \quad of which $x$ & $(n_x+1)n_yn_z$ & 385 \\
    \quad of which $y$ & $n_x(n_y+1)n_z$ & 400 \\
    \quad of which $z$ & $n_xn_y(n_z+1)$ & 420 \\
    primal cells   & $n_xn_yn_z$ & 350 \\
    dual nodes     & $n_xn_yn_z$ & 350 \\
    dual edges     & equal to the primal faces & 1205 \\
    dual faces     & equal to the primal edges & 1382 \\
    dual cells     & equal to the primal nodes & 528 \\
    \bottomrule
  \end{tabular}
\end{table}

The \emph{dual} complex is shifted by half a cell in every direction. A dual node sits at the centre of every
primal cell. A dual edge joins the two dual nodes on either side of a primal face and therefore \emph{pierces}
that face. A dual face spanned by the four dual edges around a primal edge therefore \emph{crosses} that edge. A
dual cell encloses a primal node and is bounded by the dual faces that cross the primal edges incident on that
node.

The counts in Table~\ref{tab:counts} are not an accident. For a single brick the alternating sum satisfies
\begin{equation}
  n_0 - n_1 + n_2 - n_3 = 528 - 1382 + 1205 - 350 = 1,
  \label{eq:euler}
\end{equation}
which is the Euler characteristic of a ball and holds for every grid; it is the reason the primal and dual counts
pair up as they do. In particular each dual edge threads exactly one primal face, and each dual face is pierced by
exactly one primal edge, so the pairing of primal and dual objects is a bijection. The implementation exploits
this: dual objects carry no indices of their own and are addressed by the primal object they are paired with,
which is why the code has exactly three index families -- edges, faces, cells -- and not six.

\subsection{Degrees of freedom as integrals}

Each field component is stored as its integral over the object it is naturally associated with,
\begin{align}
  \hat e_i &= \int_{L_i} \vec e\cdot\dd\vec l \quad [\mathrm{V}], &
  \hat h_f &= \int_{L_{\tilde f}} \vec h\cdot\dd\vec l \quad [\mathrm{A}], &
  \hat b_f &= \int_{A_f} \vec b\cdot\dd\vec A \quad [\mathrm{Wb}], \label{eq:dof-a}\\[2pt]
  \hat d_i &= \int_{A_{\tilde i}} \vec d\cdot\dd\vec A \quad [\mathrm{C}], &
  \hat\jmath_i &= \int_{A_{\tilde i}} \vec j\cdot\dd\vec A \quad [\mathrm{A}], &
  \hat q_c &= \int_{V_c} \rho\,\dd V \quad [\mathrm{C}]. \label{eq:dof-b}
\end{align}
The electric quantities $\mathbf{e}$ and $\mathbf{d}$ live on primal edges and the dual faces paired with them; the
magnetic quantities $\mathbf{h}$ and $\mathbf{b}$ live on the dual edges paired with primal faces and on those
primal faces themselves. The electromagnetic unknowns of the method are thus the pair $(\mathbf{e},\mathbf{h})$,
and $\mathbf{b}$, $\mathbf{d}$ are recovered from them through the constitutive relations of
Section~\ref{sec:hodge}.

\begin{proposition}[Exactness of the integral form]
\label{prop:exactness}
Let the fields satisfy Maxwell's equations in the continuum, and let a cell of the grid be given. Then Faraday's
law applied to the boundary of a primal face, Ampere's law applied to the boundary of a dual face, the divergence
theorem applied to the boundary of a primal cell and to the boundary of a dual cell are \emph{exact} relations
among the numbers \eqref{eq:dof-a}, \eqref{eq:dof-b} attached to that cell. No approximation is involved in
passing from the continuum to these relations.
\end{proposition}

\begin{proof}[Proof sketch]
Apply Stokes' theorem on a primal face: the surface integral of the curl of $\vec e$ equals the circulation of
$\vec e$ around the boundary, which by construction splits into the four edge integrals $\hat e_i$ with signs
$+1$ or $-1$. Apply the divergence theorem on a primal cell: the volume integral of the divergence of $\vec b$
equals the flux through the six bounding faces, which splits into the six face integrals $\hat b_f$. The dual
statements follow by applying the same two theorems to dual faces and dual cells.
\end{proof}

This is the property that motivates the whole method. The integral form of Maxwell's equations is not an
approximation of the differential form that becomes accurate for small cells -- it is equivalent to it, and the
two are related by the same two theorems. What a discretisation can get wrong is therefore not the differential
operator but only the constitutive law, which is what the next three subsections make precise.

\subsection{Topological operators}

The three operators are matrices of incidence. All of their entries are drawn from $\{-1,0,+1\}$; no length, area
or timestep enters any of them.

\begin{definition}[Gradient]
$\Gmat\in\{-1,0,+1\}^{n_1\times n_0}$ has one row per primal edge, with $+1$ at the node at the tail of the edge
and $-1$ at the node at its head (with a fixed orientation convention for the three edge families). Applied to a
nodal scalar potential it returns the potential difference along each edge, $\Gmat\boldsymbol\varphi$.
\end{definition}

\begin{definition}[Curl]
$\Cmat\in\{-1,0,+1\}^{n_2\times n_1}$ has one row per primal face and one nonzero entry per edge of that face,
$+1$ or $-1$ according to whether the edge is traversed with or against the face's orientation. Then
$(\Cmat\mathbf{e})_f$ is the circulation of $\mathbf{e}$ around the boundary of face $f$, so that $\Cmat\mathbf{e}$
is the discrete curl of the edge field.
\end{definition}

\begin{definition}[Divergence]
$\Smat\in\{-1,0,+1\}^{n_3\times n_2}$ has one row per primal cell and one nonzero entry per bounding face, $+1$
for the face in the positive direction along the face's axis and $-1$ for the opposite one. Then
$(\Smat\mathbf{b})_c$ is the net outward flux of $\mathbf{b}$ through the boundary of cell $c$.
\end{definition}

The dual operators are the transposes: $\tilde\Smat = \Smat^{\!\top}$ maps a cell quantity to its face fluxes, and
$\tilde\Cmat = \Cmat^{\!\top}$ maps a face quantity to a dual edge quantity. That transposition is not a coincidence
but the discrete shadow of integration by parts on the dual pair of complexes.

\begin{proposition}[The two exact identities]
\label{prop:identities}
For every orthogonal grid, uniform or not, and in exact integer arithmetic,
\begin{equation}
  \Cmat\Gmat = \mathbf{0} \quad\text{and}\quad \Smat\Cmat = \mathbf{0}.
  \label{eq:identities}
\end{equation}
\end{proposition}

\begin{proof}[Proof sketch]
$(\Cmat\Gmat\boldsymbol\varphi)_f$ is a sum of potential differences taken once around the closed loop bounding
face $f$. Every node of that loop is visited exactly once with $+1$ and once with $-1$, so the sum telescopes to
zero. For the second identity, $(\Smat\Cmat\mathbf{e})_c$ sums the circulations around the six faces of cell $c$
with outward signs; every edge of the cell belongs to exactly two of those faces and is traversed in opposite
directions by them, so every contribution appears twice with opposite signs.
\end{proof}

\begin{remark}[Metric freedom]
Because the entries of $\Cmat$, $\Gmat$ and $\Smat$ are integers, the identities \eqref{eq:identities} cannot
degrade under refinement: they are not approximately satisfied, they are satisfied. Nothing in the three matrices
refers to a length scale, so moving the nodes -- grading the mesh, refining one corner of the domain -- leaves the
operators bit-for-bit identical, provided only the counts are unchanged. This is asserted directly by comparing a
uniform and a non-uniform grid of the same shape. All geometric information of the problem is therefore collected
in one place, the Hodge stars of Section~\ref{sec:hodge}, and that separation is the structural reason a graded
grid costs nothing here.
\end{remark}

Taken together the three operators form the discrete de Rham complex
\begin{equation}
  \text{scalar potential} \xrightarrow{\ \Gmat\ } \text{primal edges} \xrightarrow{\ \Cmat\ } \text{primal faces}
  \xrightarrow{\ \Smat\ } \text{primal cells},
  \label{eq:complex}
\end{equation}
with the dual complex running the other way, $\tilde\Smat$ from cells to faces and $\tilde\Cmat$ from faces to dual
edges. Proposition~\ref{prop:identities} says that \eqref{eq:complex} is a chain complex: the image of each map
lies in the kernel of the next. Exactness -- the converse inclusion, image equal to kernel -- holds on a simply
connected domain without boundary conditions and is precisely what fails once a conductor cuts the domain; the
consequence is the gradient null space of the eigenmode section.

\subsection{The Maxwell grid equations}

With the degrees of freedom of \eqref{eq:dof-a}, \eqref{eq:dof-b} and the operators above, Maxwell's equations
read
\begin{align}
  \Cmat\,\mathbf{e} &= -\frac{\dd}{\dd t}\bigl(\Mm\,\mathbf{h}\bigr), \label{eq:faraday}\\
  \tilde\Cmat\,\mathbf{h} &= \frac{\dd}{\dd t}\bigl(\Em\,\mathbf{e}\bigr) + \mathbf{j}, \label{eq:ampere}\\
  \Smat\bigl(\Em\,\mathbf{e}\bigr) &= \mathbf{q}, \label{eq:gauss-e}\\
  \Smat\bigl(\Mm\,\mathbf{h}\bigr) &= \mathbf{0}, \label{eq:gauss-b}
\end{align}
where $\mathbf{b} = \Mm\mathbf{h}$ and $\mathbf{d} = \Em\mathbf{e}$ are the constitutive relations, and $\Mm$,
$\Em$ are the diagonal Hodge star matrices of Section~\ref{sec:hodge}. Equations \eqref{eq:faraday} and
\eqref{eq:ampere} are Faraday's and Ampere's laws on primal and dual faces; \eqref{eq:gauss-e} and
\eqref{eq:gauss-b} are the two Gauss laws on primal cells.

\begin{remark}[Where the approximation went]
Equations \eqref{eq:faraday}--\eqref{eq:gauss-b} are exact. They were obtained by applying the integral form to
the faces and cells of the grid, which Proposition~\ref{prop:exactness} permits without approximation, and no Taylor expansion has
been performed. The single approximate step of the entire method is the replacement of the constitutive laws by
the diagonal matrices $\Em$ and $\Mm$. Everything that can go wrong with this discretisation -- dispersion,
staircasing, the wrong energy at short wavelengths -- originates in that one step.
\end{remark}

The two Gauss laws need not be imposed: they are consequences of \eqref{eq:faraday}, \eqref{eq:ampere} and the
identities \eqref{eq:identities}, and are preserved in time. Differentiating \eqref{eq:gauss-b} and substituting
Faraday's law,
\begin{equation}
  \frac{\dd}{\dd t}\,\Smat\bigl(\Mm\mathbf{h}\bigr) = -\Smat\Cmat\,\mathbf{e} = \mathbf{0},
  \label{eq:gauss-preserved}
\end{equation}
so magnetic flux stays divergence free if it starts that way; the electric counterpart follows the same way from
the dual identity $\tilde\Cmat\tilde\Smat = (\Smat\Cmat)^{\!\top} = \mathbf{0}$. The implementation integrates only
\eqref{eq:faraday} and \eqref{eq:ampere} and relies on this for the rest.

\subsection{The Hodge star (material) matrices}
\label{sec:hodge}

The constitutive relations are where the metric, and with it the physics of the material, enters. Reconstruct the
field inside the cell as uniform: the component of $\vec e$ along primal edge $i$ is $\hat e_i/L_i$, and the flux
of $\vec d$ through the dual face $\tilde A_i$ perpendicular to that edge is then
$\eps_r\eps_0\,\hat e_i A_{\tilde i}/L_i$. The same argument on the magnetic side, where $\hat h_f$ is a line
integral over a dual edge and $\vec b$ is fluxed through the primal face, gives
\begin{equation}
  [\Em]_{ii} = \eps_{r,i}\,\eps_0\,\frac{A_{\tilde i}}{L_i}\ \ [\mathrm{F}],
  \qquad\qquad
  [\Mm]_{ff} = \mu_{r,f}\,\mu_0\,\frac{A_f}{L_{\tilde f}}\ \ [\mathrm{H}],
  \label{eq:stars}
\end{equation}
with all off-diagonal entries zero: the primal edge $i$ and its dual face are perpendicular by construction, and
on an orthogonal grid no other pair contributes. The ratio of the dual measure to the primal measure is the entire
content of the star; the material parameters $\eps_r$ and $\mu_r$ may be scalars, per-edge and per-face arrays, or
callables of position, in which case they are sampled at the centre of the corresponding object.

The choice of a single reconstructed component per cell is an averaging assumption. It is second order accurate on
orthogonal grids because the primal edge is centred on its dual face and the primal face is centred on its dual
edge, so the midpoint rule is exact for linear fields. The \emph{units} of \eqref{eq:stars} are the check that the
whole scheme is consistent. $\Em$ has units of capacitance and $\Mm$ of inductance, and therefore
\begin{equation}
  U_{\mathrm{e}} = \tfrac12\,\mathbf{e}\cdot(\Em\mathbf{e})\ \ [\mathrm{J}], \qquad
  U_{\mathrm{m}} = \tfrac12\,\mathbf{h}\cdot(\Mm\mathbf{h})\ \ [\mathrm{J}], \qquad
  P = \mathbf{e}\cdot\mathbf{j}\ \ [\mathrm{W}],
  \label{eq:energies}
\end{equation}
which is the normalisation the energy analysis of the section on energy conservation and stability relies on.

Two consequences deserve emphasis. First, because the stars are diagonal, a product with $\Em$ or $\Mm$ is an
elementwise scaling and never a linear solve. A time step therefore costs two sparse matrix-vector products by
$\Cmat$ and $\tilde\Cmat$ plus four elementwise scalings, and nothing else; this is what makes an explicit scheme
viable in three dimensions.

Second, the star is \emph{not} length invariant, unlike the topological operators. Since $A_{\tilde i}/L_i$ has
dimension of length, scaling every cell by a factor $s$ multiplies every entry of $\Em$ and of $\Mm$ by exactly
$s$ -- the test suite asserts precisely this for $s=2$. As a grid is refined the stars shrink linearly with the
cell size, which is the origin of the linear-in-$h$ CFL bound of the section on energy conservation and stability.

\begin{remark}[Graded and non-orthogonal grids]
On an orthogonal grid with non-uniform spacing the stars remain diagonal -- only their entries change -- and this
is why grading costs nothing structurally. On a non-orthogonal (skew) grid the reconstruction inside a cell mixes
field components, so the stars become sparse, symmetric positive definite, banded matrices. The topology
operators, and hence Proposition~\ref{prop:exactness}, are untouched; what changes is that the explicit step now needs an
approximate inverse of a sparse matrix. That case is outside the present implementation, and it is the reason the
material layer is kept strictly separate from the topological operators in the code.
\end{remark}

\subsection{Boundary conditions and sources}

A perfect electric conductor prescribes vanishing tangential electric field on its surface. In the discrete model
that is the statement that the edge voltages of all edges lying in the wall vanish identically,
\begin{equation}
  \hat e_i = 0 \quad\text{for every edge } i \text{ tangential to a PEC wall}.
  \label{eq:pec}
\end{equation}
The implementation imposes \eqref{eq:pec} as a constraint on the unknown space rather than as an operation on the
solution: the unknown vector is restricted to the free edges and the corresponding columns of $\Cmat$ never
participate. This matters for exactness. A zero column contributes nothing to $\Cmat\mathbf{e}$, so the restricted
update is an exact leapfrog on a smaller complex and inherits the invariant of the section on energy conservation and stability bit for bit,
whereas the alternative of stepping everywhere and then overwriting the wall values with zero would silently
inject or remove the energy belonging to those edges. On the $10\times7\times5$ cavity with all six walls
conducting, 620 of the 1382 edge voltages are pinned this way, leaving 762 unknowns.

A perfect magnetic conductor is the dual condition: the tangential magnetic field vanishes, which in these
variables means that the magnetic flux through each wall face vanishes, $\hat b_f = 0$ for faces in the wall. PEC
therefore constrains the primal edge space and PMC the primal face space, and the two conditions are the two ends
of the duality that runs through the whole construction.

The restriction to free edges does, however, cost something, and it is worth naming it here. Proposition~\ref{prop:exactness}
concerns the full complex; after deleting the PEC unknowns the restricted curl no longer annihilates the
restricted gradient by the same integer identity, and the surviving gradient fields reappear as a null space of
the discrete curl--curl operator. For a box with all six walls conducting that null space has dimension
$(n_x-1)(n_y-1)(n_z-1)$ -- 216 on the coarse cavity grid and 702 on the fine one. The eigenmode section shows how
that cluster is identified and removed.

Sources enter in two ways. A \emph{soft} source is a current density carried by the dual faces and added to the
right-hand side of \eqref{eq:ampere}; because $\mathbf{e}\cdot\mathbf{j}$ has units of power, by \eqref{eq:energies}
it injects
\begin{equation}
  P = \mathbf{e}\cdot\mathbf{j}\ \ [\mathrm{W}],
  \label{eq:power}
\end{equation}
which is the only channel by which energy enters the system and is used verbatim in the section on energy conservation and stability. A
\emph{hard} source, by contrast, prescribes the value of $\hat e_i$ on a chosen edge and overrides the update
there; the example drives the cavity with a Ricker waveform imposed this way. Losses would enter as the
conduction current $\mathbf{j} = \Ms\mathbf{e}$; $\Ms$ is assembled with the right units but is not yet coupled
into the step, which is why every statement about energy conservation in this report is a statement about a
lossless problem.

\subsection{Relation to Yee's scheme and to finite elements}

Choose a uniform orthogonal grid and identify the primal edge voltages with the components of $\vec e$ at the
primal edge midpoints and the dual edge magnetic voltages with the components of $\vec h$ at the primal face
centres. Eliminating $\mathbf{b}$ and $\mathbf{d}$ through the diagonal stars, equations \eqref{eq:faraday} and
\eqref{eq:ampere} become exactly the staggered update of Yee \cite{yee1966}. FIT is thus not a competitor of the
finite-difference time-domain method but its structural generalisation: the staggering is the same, while the
metric has been isolated into an explicit object, which is what allows non-uniform spacing without touching a
single operator. On the uniform grid the classical results on numerical dispersion, the CFL condition and the
staircase error of curved boundaries \cite{taflove2005} apply verbatim.

The same de Rham complex underlies edge-element finite elements \cite{bossavit1998,jin2014}, where Whitney one-forms
carry $\vec e$ and two-forms carry $\vec b$; the incidences are the same integers. The difference is entirely in
the Hodge star. A finite-element assembly builds a sparse mass matrix from element integrals, which represents
curved and anisotropic geometry accurately but must be solved at every step. FIT replaces it with a diagonal,
measure-based, low-order approximation, which is why an explicit step exists at all. Bossavit's complementarity
view puts the two families on one footing; for the problems this project targets the trade of geometric fidelity
against cost per unknown falls on the FIT side, at the price of a structured grid and of a staircase at curved
surfaces.

\begin{remark}[What is exact in the scheme, and what is not]
\label{rem:exact}
The word ``exact'' has been used three times above with three different objects, and it is worth separating them
once, because the whole credibility of the numbers in this report rests on the separation. The topological
identities are exact in the strongest sense: $\Smat\Cmat=0$ and its companions are statements about integer
incidence matrices, they hold for any material distribution, on any grid, at any time step, and they are what
Proposition~\ref{prop:identities} and \eqref{eq:identities} record. The constitutive relation is not exact: the
Hodge star matrices $\Em$, $\Mm$ of \eqref{eq:stars} replace a continuous field by a diagonal, measure-based,
low-order average, and every approximation in space is in that replacement and nowhere else
(Section~\ref{sec:hodge}). The time step is not exact either: the leapfrog recursion
\eqref{eq:leapfrog-h}--\eqref{eq:leapfrog-e} conserves the mixed-state quadratic form \eqref{eq:invariant} to
rounding, which is an exact statement about the discrete system, but the scheme is only conditionally stable
\eqref{eq:cfl}, at a step set by the smallest cell (Table~\ref{tab:cfl}), and stepping inside the bound does not
remove the dispersion of \eqref{eq:dispersion}, which is why a tracked wavefront travels at $0.904\,c_0$
(Section~\ref{sec:dispersion}). Nothing in the report therefore claims that the finite integration technique is
exact. What can be claimed is narrower and is measured: the conserved quantity is exact, the incidence relations
are exact, and the accuracy of any field, frequency or impedance quoted here is the accuracy of the stars at the
resolution used, which for the closed cavity is about one per cent on a coarse grid with second-order convergence
(Section~\ref{sec:results}). An error in that accounting, whether it is attributed to the grid when it belongs to
the boundary or the other way round, is the one mistake this scheme makes easy, and the measurements of
Section~\ref{sec:worked} are reported per record for that reason.
\end{remark}

